\chapter{Introdução}

A robótica e a automação industrial têm passado por transformações substanciais nas últimas décadas, deixando de ser exclusividade de grandes corporações para se integrarem a ambientes acadêmicos e laboratórios de pesquisa. Impulsionadas pela democratização do hardware e por abordagens \textit{open-source}, essas tecnologias possibilitaram a fabricação de braços robóticos acessíveis com componentes padronizados e impressão 3D \cite{silva2020}. Desse modo, a construção de protótipos funcionais tornou-se viável para explorar arquiteturas de controle e instrumentação sem os custos proibitivos do setor industrial, representando um pilar essencial para a expansão da robótica educacional e experimental.

Nesse contexto de conectividade crescente, a inserção do paradigma de Internet das Coisas (IoT) em sistemas robóticos introduziu novas possibilidades de monitoramento, controle remoto e gerenciamento distribuído de dados. Conforme \citeonline{siciliano2009}, arquiteturas conectadas permitem que periféricos físicos interajam com interfaces web, estabelecendo ecossistemas em rede alinhados ao conceito da Indústria 4.0. Todavia, a implementação dessas arquiteturas em plataformas de baixo custo esbarra em desafios de processamento embarcado, sincronização de atuadores e estabilidade da comunicação em tempo real.

Diante desse cenário, o presente trabalho concentra-se no desenvolvimento de um protótipo de braço robótico de baixo custo, baseado no modelo \textit{open-source} EB-15, combinando acionamento motorizado, controle eletrônico distribuído e funcionalidades de rede. A motivação surge da necessidade de criar uma infraestrutura experimental acessível, capaz de executar tarefas básicas de manipulação e de hospedar uma interface de operação local inspirada em sistemas supervisórios industriais. Sendo assim, o capítulo detalha a fundamentação da pesquisa, apresentando o problema central, os objetivos e a metodologia proposta.

\section{Problema de Pesquisa}

Em uma aplicação contemporânea de braço robótico, questiona-se a viabilidade de um protótipo de baixo custo com integração IoT executar funções básicas de manipulação e apresentar desempenho funcional comparável, em critérios selecionados, a um robô industrial de referência. Essa interrogação impõe a necessidade de investigar não apenas a capacidade construtiva do sistema, mas também a competência do \textit{firmware} embarcado em coordenar múltiplos eixos motorizados enquanto mantém serviços ativos na rede local. Portanto, a viabilidade técnica do protótipo constitui o eixo central desta pesquisa.

De forma complementar, indaga-se quais são as principais limitações e potencialidades desse tipo de sistema quanto à precisão posicional, à repetibilidade de trajetórias, à estabilidade dinâmica e à capacidade de execução de movimentos coordenados. Tal abordagem direciona a análise para a quantificação experimental do comportamento do equipamento, confrontando métricas objetivas com aquelas obtidas em plataformas industriais consolidadas. Por conseguinte, a convergência dessas duas vertentes investigativas orienta tanto o projeto de engenharia de controle quanto a metodologia de testes adotada.

\section{Objetivos}

\subsection{Objetivo Geral}

O objetivo principal deste trabalho consiste em desenvolver e validar um protótipo de braço robótico de baixo custo com integração IoT para execução de funções básicas de manipulação. Em complemento, busca-se investigar a viabilidade técnica de uma solução baseada em conectividade distribuída para aplicações experimentais em células de automação acadêmicas.

\subsection{Objetivos Específicos}

Para alcançar a meta geral estabelecida, foram delineados objetivos específicos que servirão de roteiro para o desenvolvimento do projeto. Primeiramente, pretende-se projetar a estrutura mecânica e a eletrônica embarcada do braço robótico, definindo motores, redutores e sensores adequados à proposta de baixo custo. Na sequência, busca-se implementar o sistema de controle de baixo nível e a comunicação IoT do protótipo por meio de uma arquitetura baseada em microcontroladores, além de desenvolver uma interface web de operação que viabilize o comando manual de eixos, o monitoramento de status e a parametrização do braço.

Ademais, objetiva-se executar testes práticos de movimentação básica, englobando posicionamento espacial, deslocamentos lineares e repetição contínua de rotas definidas na metodologia. Soma-se a isso a comparação empírica dos resultados métricos do protótipo com dados fornecidos por um robô industrial consolidado, considerando como dimensões de avaliação o tempo de resposta, a estabilidade e a corretude posicional. Por fim, propõe-se explorar a implementação de lógicas intermediárias, como rotinas programadas de \textit{pick and place}, avaliando possíveis distorções frente à plataforma de referência.

\section{Justificativa}

A construção de equipamentos robóticos funcionais em ambiente universitário é uma ação altamente custosa quando pautada em importação e na aquisição de ferramentas proprietárias. Com efeito, explorar as capacidades reais de impressões 3D combinadas com controladores acessíveis como Arduino e ESP apresenta-se como um avanço para a formação em engenharia e áreas afins, visto que viabiliza bancadas didáticas de automação \cite{siciliano2009}. Dessa forma, esse mapeamento reduz a carência de equipamentos e impulsiona pesquisas independentes focadas em otimização de \textit{drivers} lógicos de passo, cinemática embarcada e interfaces homem-máquina (IHM) leves.

Sob o mesmo ponto de vista, o advento da integração contínua de dispositivos físicos na rede local (conceito IoT) requer estruturas de testes que unifiquem acionamentos robustos de corrente contínua à transmissão ágil de pacotes de dados. Não obstante, os laboratórios nem sempre dispõem de braços articulados dotados de interfaces abertas para experimentação em alto nível, limitando estudos sobre arquiteturas híbridas e latência de comando de motores operados remotamente. Por conseguinte, este trabalho justifica-se por fornecer não apenas o relato de desenvolvimento de hardware, mas um caso de uso concreto de controle distribuído em tempo real acessível à comunidade pesquisadora.

\section{Metodologia Aplicada}

O desenvolvimento deste trabalho segue uma abordagem metodológica estruturada em fases sequenciais e interdependentes. Primeiramente, realiza-se o levantamento bibliográfico e a fundamentação teórica, revisando conceitos de projeto mecânico de manipuladores robóticos, atuadores, protocolos de integração voltados para Internet das Coisas, arquitetura de microcontroladores modernos e as equações matemáticas de cinemática aplicáveis ao sistema.

Em seguida, detalha-se a fase de implementação prática do protótipo, abrangendo a especificação dos componentes físicos do modelo adotado (EB-15), o esquemático das conexões elétricas estabelecidas, a arquitetura dual de controle utilizada (processadores Mestre e Escravo) e as escolhas de projeto da interface web de operação desenvolvida.

Posteriormente, apresenta-se o protocolo experimental de testes e a análise dos resultados obtidos. Nessa etapa, abordam-se comparativamente o desempenho do protótipo em tarefas contínuas (repetição guiada, movimentos coordenados) e as métricas coletadas frente ao robô industrial selecionado como referência. Por fim, discutem-se as considerações finais, analisando o grau de atendimento dos objetivos estipulados e fornecendo perspectivas sobre possíveis melhorias arquiteturais e expansões do \textit{software} implementado.
